\paragraph{Minorant profiles and disjoint owned choice.}
Disjoint union combines activation preorders directly. For owned choice,
the remaining singleton stratum is determined by a second piece of data.
For a core operation $C$ on finite nonempty $V$, define its
\emph{minorant profile envelope} at owner $s$ by
\[
 M_s^C(A)=\bigvee_{G\text{ core},\ G\subseteq C}\Phi_s(G)(A).
\]
This whole profile is attained by one minorant: there are finitely many
core operations at this arity, their minorants are closed under every
owned choice, and a $\star_s$-fold disjoins their $s$-profiles.
The envelope is monotone, is zero at $\varnothing$, and is one at $V$.

\begin{theorem}[Minorant profiles under disjoint composition]
\label{core:thm:disjoint-minorant-profiles}
Let $K\geq3$, let $C_i$ be arbitrary core operations on disjoint finite
nonempty $V_i$, and write $A_i=A\cap V_i$. Then
\[
 \begin{aligned}
 M_s^{C_1+C_2}(A)&=M_s^{C_1}(A_1)\lor M_s^{C_2}(A_2),\\
 M_r^{C_1\star_r C_2}(A)&=M_r^{C_1}(A_1)\lor M_r^{C_2}(A_2),\\
 M_s^{C_1\star_r C_2}(A)&=M_s^{C_1}(A_1)\land M_s^{C_2}(A_2)
       \qquad(s\ne r).
 \end{aligned}
\]
\end{theorem}
\begin{proof}
For the union upper bound, let $G\subseteq C_1+C_2$ and
$\Phi_s(G)(A_1\cup A_2)=1$. If $A_2=\varnothing$, neutral deletion
of $V_2$ already gives a minorant of $C_1$ positive at $A_1$.
Otherwise neutrally delete $V_2\setminus A_2$, identify the remaining
$A_2$ inputs with $Y$, and fix $Y=\{s\}$. The nonzero $s$-profile
survives the first deletion exactly. If the fixed cofactor is not
constantly $\{s\}$, fixed-menu erasure gives a core operation positive
at $A_1$ and bounded by
$\partial_{V_2}G\subseteq C_1$.
If it is constantly $\{s\}$, the exact profile preservation gives
$\Phi_s(G)(A_2)=1$ already before deletion. Neutral deletion of $V_1$
then gives a minorant of $C_2$ positive at $A_2$.
This proves the disjunction. Conversely, combine minorants attaining
the two child envelopes by $\star_s$. They lie below $C_1+C_2$, and
their $s$-profiles disjoin.

Put $C=C_1\star_r C_2$. A minorant $G\subseteq C$ satisfies
\[
 \partial_{V_2}G\subseteq C_1,\qquad
 \partial_{V_1}G\subseteq C_2,\qquad
 \Phi_r(G)(V_1)=\Phi_r(G)(V_2)=1.
\]
Indeed $\partial_{V_2}C=C_1$ and $\partial_{V_1}C=C_2$:
at each winner's neutral menu, the other child is that same ambiguous
menu by idempotence. Fixing either whole block to $r$ makes $C=\{r\}$.
For $s\ne r$, suppose $\Phi_s(G)(A_1\cup A_2)=1$.
Both $A_i$ are nonempty, by compatibility with the two displayed
$r$-coalitions. Perform the same partial deletion and fixed-$s$
erasure as above. Its constant-$s$ alternative is now impossible:
the preserved $r$-profile at $V_1$ still makes the cofactor equal
$\{r\}$ there. Thus the erasure is a minorant of $C_1$ positive
at $A_1$. Interchanging the blocks gives the analogous minorant of
$C_2$. This proves the foreign-owner conjunction.
This extraction argument uses only the two residual bounds and the
two $r$-profile bits displayed above, not the full bound $G\subseteq C$.

For owner $r$, the union formula is an upper bound since
$C\subseteq C_1+C_2$. For every owner, combining envelope-attaining
child minorants by $\star_r$ stays below $C$ and attains the stated
disjunction or conjunction. This completes all three equalities.
\end{proof}

\begin{theorem}[Activation under disjoint owned choice]
\label{core:thm:disjoint-owned-activation}
Under the same hypotheses, put $C=C_1\star_r C_2$.
For proper colorings $p,E$ on $V_1\cup V_2$, its principal activation
predicate is
\[
 \Theta_{C,p}(E)=
 \begin{cases}
  \Theta_{C_1,p_1}(E_1)\,\Theta_{C_2,p_2}(E_2),
       &\Phi_r(C)(p_r)=0,\\[2pt]
  \displaystyle\bigwedge_{s\ne r}
       \neg\bigl(M_s^{C_1}(E_s\cap V_1)\,
                 M_s^{C_2}(E_s\cap V_2)\bigr),
       &\Phi_r(C)(p_r)=1.
 \end{cases}
\]
\end{theorem}
\begin{proof}
The first case is \cref{core:cor:owned-choice-activation,core:thm:disjoint-activation-product}.
For the second, singleton-fiber constancy lets us take $p$ all-$r$.
For each $s\ne r$, choose child minorants $G_s^i$ attaining
$M_s^{C_i}$. The helper $G_s^1\star_r G_s^2$ lies below $C$ and
has $s$-profile equal to the displayed product.
Attach these helpers to $X$ at their owners $s$, then adjoin $C$.
The resulting interval operation has exactly the stated box and is
positive at $p$.

For minimality take a principal one-$X$ spine $S\subseteq C+X$
positive at $p$. It has no $r$-owned path gate, since every passive
helper is $\{r\}$ on the all-$r$ diagonal.
Every helper $H$ satisfies
$\Phi_r(H)(V_1)=\Phi_r(H)(V_2)=1$:
when $X=\{r\}$ and either whole passive block is fixed to $r$,
the bound makes the spine the $r$ singleton, and its foreign path
gates conjoin all these helper profile bits.

Give every input in $V_2$ an ambiguous menu. No helper can be a
foreign singleton, since its foreign winning coalition would be
disjoint from the positive $r$-coalition $V_2$.
With $X=I\setminus\{k\}$, every path gate consequently takes union.
Whenever $C$ omits $k$, the spine bound forces every helper to omit
$k$ too. Thus $H\subseteq C$ on these rows, and winnerwise neutral
deletion gives $\partial_{V_2}H\subseteq C_1$.
The analogous bound holds after deleting $V_1$.
Apply the preceding proof's foreign-profile extraction to $H$.
It bounds each owned helper profile by the corresponding product
$M_s^{C_1}M_s^{C_2}$. The spine's $s$-mask is the disjunction of
these helper profiles, and its $r$-mask is zero.
Its box therefore contains the box already attained, proving equality.
\end{proof}

Thus activation and minorant profile envelopes together admit exact
composition rules for disjoint union and disjoint owned choice.
The ordinary singleton profiles used to select the branch obey their
usual gate rules: owned choice disjoins its owner's bits and conjoins
the others, while union conjoins every owner's bits.
The envelopes concern bounded core operations, rather than just the
singleton profiles of the baseline. Explicit starting cases below give
closed tests for disjoint compositions without searching for minorants.

\begin{proposition}[Minorant envelopes see more than singleton profiles]
\label{core:prop:minorant-profile-data}
For every $K\geq3$ there are three-input core operations $T,D$ with
the same singleton profiles but different minorant profile envelopes.
Specifically, for distinct owners $a,b,c$, put
\[
 T=(u\star_a x)\star_b(x\star_c v)
\]
and let $D$ be its greatest profile completion. Writing variable names
for coalition membership bits, one has
\[
 M_b^T=x(u\lor v),\qquad M_b^D=x.
\]
\end{proposition}
\begin{proof}
The profiles of $T$ are
$\Phi_a=xv$, $\Phi_c=xu$, $\Phi_b=x(u\lor v)$, and $xuv$ at all
other owners. Any minorant preserves the winning $a$-coalition
$\{x,v\}$ and $c$-coalition $\{x,u\}$, so every winning
$b$-coalition must contain $x$ by rainbow compatibility.
It cannot be $\{x\}$: at $w_u=\{a\},w_x=\{b\},w_v=\{c\}$,
$T(w)=\{a,c\}$, whereas that positive profile bit would force a
minorant to return $\{b\}$. Hence $M_b^T=\Phi_b(T)$.
Replacing the $b$-profile by $x$ and retaining all other profiles
gives an admissible tuple: every winning coalition contains $x$.
Its core realizer lies below the profile maximum $D$, because it
preserves every singleton value of $D$. Thus $M_b^D\geq x$, and
the same two-coalition compatibility bound gives equality.
\end{proof}

Disjointness is also essential to the composition laws themselves.
For distinct $a,b,c$, derived union gives
$(x\star_a y)\star_c(x\star_b y)=x+y$.
Each child has $c$-envelope $xy$, while the result has $c$-envelope
$x\lor y$ by the union law on its two inputs. Identifying child inputs
can therefore create new envelope bits that the disjoint rule does not supply.
